%% Section 1: Introduction
%% Paper: Littlewood-Richardson Coefficients: From Puzzles to the Integrability Hierarchy

\section{Introduction}\label{sec:introduction}

The Littlewood-Richardson coefficients $c^\lambda_{\mu\nu}$ are among the most
studied numbers in algebraic combinatorics.  They appear as the structure
constants of the ring of symmetric functions,
\begin{equation}\label{eq:lr-def}
  s_\mu \cdot s_\nu \;=\; \sum_\lambda c^\lambda_{\mu\nu}\, s_\lambda,
\end{equation}
as tensor product multiplicities of $GL_n$ representations, and as intersection
numbers in the Schubert calculus of Grassmannians.  What is remarkable is not
that they can be computed---they can, in many ways---but that so many
\emph{independent} combinatorial interpretations coexist.  The
Littlewood-Richardson rule counts skew tableaux satisfying a lattice word
condition.  The Knutson-Tao-Woodward puzzle model~\cite{knutson-tao-woodward-2004}
tiles a triangle with three species of unit rhombi.  Purbhoo's
mosaics~\cite{purbhoo-2007} give a hybrid between the two.  Pipe dreams and
bumpless pipe dreams compute the same numbers via reduced decompositions in the
symmetric group.  Gaetz and Gao's interlacing triangular
arrays~\cite{gaetz-gao-2024} provide yet another model, discovered from
the geometry of Coxeter groups.  Each of these was found independently, often
by different communities---algebraists, combinatorialists, geometers---and the
bijections connecting them, while known, have the character of ad hoc
constructions.  One is left with a question that is more structural than
computational: \emph{why do all of these models exist?}

\subsection*{The integrable perspective}

The first hint of a deeper explanation came from Zinn-Justin's
observation~\cite{zinn-justin-2008} that KTW puzzle pieces \emph{are} the
Boltzmann weights of an exactly solvable lattice model.  More precisely, the
three puzzle tiles carry weights that satisfy the \emph{Yang-Baxter equation}
(YBE), the fundamental integrability condition for two-dimensional statistical
mechanics.  This is not a metaphor.  The row transfer matrix of the lattice
model literally processes the puzzle triangle row by row, and its partition
function evaluates to $c^\lambda_{\mu\nu}$.  The YBE guarantees that the
transfer matrices for different rows commute, which is why the boundary data
(the partitions $\lambda$, $\mu$, $\nu$) determine the count regardless of the
internal tiling---exactly as they should, since the LR coefficient depends on
the partitions alone.

Bump and Naprienko~\cite{bump-naprienko-2025} elevated this observation to a
structural principle.  They introduced the \emph{Yang-Baxter groupoid}: its
objects are combinatorial models (each specified by a choice of boundary
conditions and local vertex weights), and its morphisms are R-matrix
transformations arising from the YBE.  Two models connected by a morphism
compute the same structure constants.  Within this framework, the ``mysterious''
bijections between puzzles, tableaux, and pipe dreams become groupoid morphisms
---canonical consequences of the Yang-Baxter equation rather than feats of
combinatorial ingenuity.  At a still deeper level, Garbali and
Gunna~\cite{garbali-gunna-2024} showed that all vertex models in this story
are representations of the Feigin-Odesskii \emph{shuffle algebra}, providing a
single algebraic substrate from which the entire zoo of models can be derived.

\subsection*{The categorical phase diagram}

The Yang-Baxter groupoid explains the horizontal proliferation of models: many
combinatorial objects, all computing the same number, related by R-matrix moves.
But there is a second, orthogonal direction of variation that the groupoid alone
does not capture.  The LR coefficients are the \emph{classical} case---they
govern the cohomology ring of the Grassmannian.  One can deform to K-theory
(Grothendieck polynomials), to quantum K-theory, or to spin models
(Hall-Littlewood and $q$-Whittaker polynomials).  Each deformation has its own
integrable lattice model, its own R-matrix, and its own family of structure
constants---yet all satisfy the full Yang-Baxter
equation~\cite{wheeler-zinn-justin-2016, gunna-wheeler-zj-2025}.  The
deformation axis is orthogonal to the groupoid axis: it changes \emph{which}
structure constants are computed, not the mechanism by which models proliferate.

The central contribution of this paper is a \emph{categorical phase diagram}
that organises both directions simultaneously.  The horizontal axis parameterises
the deformation type (classical, K-theoretic, quantum K-theoretic, spin).  The
vertical axis parameterises the \emph{integrability level}---the categorical
structure that governs the symmetry of the model.  We identify three levels,
governed by the Hecke algebra chain $H_q \to H_0 \to \mathbb{C}[S_n]$:

\begin{enumerate}
\item \textbf{Braided monoidal} (generic $q$).  The Hecke generators satisfy
  $(T_i - q)(T_i + 1) = 0$, and the full Yang-Baxter equation holds.
  Symmetry is governed by the braid group $B_n$.  This is the world of
  Zinn-Justin's lattice models, Bump-Naprienko's groupoid, and the rich zoo
  of combinatorial models related by R-matrix moves.

\item \textbf{Coboundary monoidal} ($q = 0$).  At $q = 0$ the Hecke relation
  becomes $T_i(T_i + 1) = 0$; equivalently, the Demazure operators
  $\pi_i = T_i + 1$ satisfy $\pi_i^2 = \pi_i$, and the braiding
  degenerates.  What survives is a \emph{coboundary} structure: the commutor
  (the categorical analogue of the swap map) still exists, but it no longer
  satisfies the braid relation---only the weaker \emph{cactus group} relations
  of Henriques and Kamnitzer~\cite{henriques-kamnitzer-2006}.  This is the
  world of column-level integrability and the ``quasi-solvable'' models of
  Buciumas and Scrimshaw~\cite{buciumas-scrimshaw}.

\item \textbf{Symmetric monoidal} (crystal limit).  At the crystal level,
  $q$ is specialised away entirely, and only the symmetric group $S_n$ acts.
  The commutor becomes the trivial swap.  Crystal bases~\cite{kashiwara-1990}
  provide a purely combinatorial framework: Kashiwara operators, crystal graphs,
  tensor product rules.
\end{enumerate}

Each downward transition forgets something precise.  To state what, we recall
the \emph{Schur-Weyl decomposition}: for the defining representation $V$ of
$GL_n$, the tensor power decomposes as
\begin{equation}\label{eq:schur-weyl}
  V^{\otimes k} \;\cong\; \bigoplus_{\lambda \vdash k} V_\lambda \otimes S_\lambda,
\end{equation}
where $V_\lambda$ is the irreducible $GL_n$-module of highest weight $\lambda$
and $S_\lambda$ is the Specht module carrying the action of $S_k$.  The two
tensor factors---the ``physical'' space $V_\lambda$ and the ``multiplicity''
space $S_\lambda$---play sharply different roles in the hierarchy.

The key result (the \emph{Multiplicity Bundle Theorem}, proved in
Section~\ref{sec:hierarchy}) states that the coboundary commutor factorises as
\begin{equation}\label{eq:commutor}
  \sigma_{[i,j]} \;=\; \mathrm{id}_{V_\lambda} \otimes \rho_{S_\lambda}(w_0),
\end{equation}
where $w_0$ is the longest element of $S_{j-i+1}$ and $\rho_{S_\lambda}$
denotes the Specht module representation.  The commutor lives entirely in the
multiplicity space---it acts as the identity on the ``physical'' irreducible
$V_\lambda$ and as the longest-element involution on the Specht module
$S_\lambda$.  The crystal functor, which retains only the combinatorial skeleton
of the representation, kills the multiplicity space: since the tensor square
$B(\omega_1)^{\otimes 2}$ of the crystal of the defining representation is
multiplicity-free (a property special to the defining representation; the
general statement requires the Henriques-Kamnitzer crystal commutor theory),
every crystal automorphism is forced to be the identity by rigidity.  The
commutor becomes invisible.

This gives concrete mathematical content to each transition in the hierarchy.
The passage from braided to coboundary forgets the \emph{spectral parameter}:
the R-matrix $R(u)$, which depends on a continuous parameter $u$ and satisfies
the full parametric YBE, degenerates to the commutor $\sigma$, which is a
single fixed morphism satisfying only the cactus relations.  The passage from
coboundary to symmetric forgets the \emph{multiplicity bundle}: the Specht
module factor in~\eqref{eq:schur-weyl}, which carries the nontrivial action
of $\sigma$, is killed by the crystal functor, leaving only the symmetric
swap.

A crucial subtlety, clarified by Kamnitzer and
Tingley~\cite{kamnitzer-tingley-2007}: the coboundary structure is already
present at generic $q$, constructed from Drinfeld's unitarised R-matrix.
At $q = 0$, the braiding dies but the coboundary commutor survives.
The transition from braided to coboundary is not a creation but an
\emph{extraction}---the coboundary structure was always there, hidden inside
the braiding, and the $q = 0$ specialisation strips away the braided
superstructure to reveal it.  Like a zipper pulled open, the debraiding
separates two interlocking layers that were always distinct but moved in
concert.

\subsection*{Plan of the paper}

Section~\ref{sec:classical} reviews the classical theory of LR coefficients:
their definition via symmetric functions, tableaux, and Schubert calculus, and
the puzzle and Hopf algebra perspectives that set the stage for integrability.
Section~\ref{sec:lattice} introduces the integrable lattice model framework of
Zinn-Justin, presenting the Yang-Baxter equation and the transfer matrix
formalism.  Section~\ref{sec:groupoid} develops the Yang-Baxter groupoid of
Bump and Naprienko, explaining how classical bijections between combinatorial
models arise as groupoid morphisms.  Section~\ref{sec:hierarchy}, the heart of
the paper, constructs the categorical phase diagram and proves the Multiplicity
Bundle Theorem.  Section~\ref{sec:shuffle} briefly discusses the shuffle algebra
substrate underlying the entire framework.
Section~\ref{sec:open} collects open questions.

The picture that emerges is, we hope, both clarifying and aesthetically
satisfying.  The many combinatorial models for LR coefficients are not
historical accidents.  They are shadows of a single algebraic structure---the
integrable R-matrix and its degenerations---cast from different angles onto the
combinatorial plane.  The categorical phase diagram tells us which angles are
available, what is preserved at each level, and what is lost in each projection.
The coefficients themselves, those nonnegative integers governing the
multiplication of Schur functions, sit at the centre like a fixed point: the one
invariant that every shadow faithfully records.
